\section{Introduction}
\label{sec:introduction}

The high-row fixed-shift route isolates a literal determinant
coefficient
\[
 A_X(n)=A_{C_X}(n),\qquad
 Z_X=\sum_nA_X(n),\qquad
 D_X=\sum_n|A_X(n)|^2.
\]
TPC-81 and TPC-82 record \(D_X\) and \(Z_X\) as the two remaining
arithmetic statistics in the conditional flatness transfer
\citep{WangTPC81,WangTPC82}.  TPC-83 reduces natural determinant
energy to survival of the actual post-bin \(\ell^1\) mass
\citep{WangTPC83}.  TPC-85 then audits several operations that do not
manufacture cancellation: Fourier centering, determinant-fiber
conditional expectation, and content inversion
\citep{WangTPC85}.

There is, however, a subtle correction to a common exclusion.
The physical TPC-32 Fourier modulus is chosen large enough that the
determinant support embeds without wrap, and the literal determinant
bin \(n=0\) vanishes.  Fourier inversion at that structural zero
expresses \(Z_X\) as a sum of all nonzero Fourier modes.  Consequently
a sufficiently coherent theorem for those modes can control \(Z_X\).
What fails is the unproved passage from generic, mean-square, or
all-but-exceptional frequency information to that particular signed
sum.

This paper makes the distinction exact.  Its main results are:
\begin{enumerate}[label=(\roman*),leftmargin=2.2em]
\item every zero coordinate gives an exact nonzero-frequency
reconstruction of the zero mode;
\item among kernels obtained by averaging structural zeros, the
unique minimum-\(\ell^2\) kernel averages them uniformly;
\item the optimal kernel and Parseval recover the sharp support
inequality \(|Z|^2\le |S|D\), with no loss;
\item positive-density zero-arc padding produces fixed-order smooth
filters with bounded Wiener norm and polynomial frequency decay;
\item the resulting tail estimate reduces the literal zero-mode
problem to a mesoscopic low-frequency window;
\item exceptional frequencies must be measured by their mass under
the reconstruction multiplier, not merely by cardinality;
\item total nonzero-frequency energy alone gives no improvement over
the sharp support bound; and
\item the literal TPC-32 packet satisfies all finite hypotheses, but none of
the finite identities supplies new arithmetic cancellation.
\end{enumerate}

The point is not to relabel the zero-mode problem as a Fourier
problem.  The exact identity instead says what a valid Fourier input
must estimate.  It also prevents two opposite errors: declaring that
nonzero frequencies are logically irrelevant, and declaring that a
generic dispersion estimate automatically controls their coherent
sum.

\subsection{Audit levels}

We use the established three-level terminology.  An \(\Lzero\)
statement is finite algebra or an abstract countermodel.  An
\(\Lone\) statement attaches that algebra to the literal
fixed-\(h_0\) carrier with complete native aggregation, actual support,
and global normalization.  An \(\Ltwo\) statement is a new
growing-\(X\) arithmetic estimate for that carrier.  The main Fourier
theorems below are \(\Lzero\); \cref{sec:physical-interface} is
\(\Lone\); the missing coherent M\"obius estimate is \(\Ltwo\).
